% Imporditud: tools/import_eksperimendid.py
% Allikas: Eesti füüsikaolümpiaadi eksperimendi ülesannete
%   kogu (Rael Kalda, 2023), ülesanne Ü156, lahendus L156.
\setAuthor{Jaan Kalda}
\setRound{lõppvoor}
\setYear{2020}
\setNumber{G E2}
\setDifficulty{4}
\setTopic{elektriahelad}
\setVahendid{must kast, millel on neli väljundjuhet (valge, roheline, sinine ja punane), oommeeter. (kokku 14 p)}
\setPunktid{14}
\setAllikas{Eksperimendi kogu Ü156, lahendus lk 117}
\setLahendusseis{toores}

\prob{Must kast}
Must kast sisaldab nelja takistit, mis on ühendatud teadmatul viisil nelja väljundjuhtme vahele. Juhtmete otsi saab ühendada oommeetri külge ja vajaduse korral ka lühistada omavahel. (a) Määrake kindlaks musta kasti sees olev elektriskeem (joonista see skeem, tähista skeemil juhtmete värvid ja põhjenda mõõtmistulemustega). (7 p) (b) Leidke kõigi takistite väärtused. (7 p)

\hint

\solu
Märkus: tegelikult poleks traati vaja, piisaks kuuest mõõtmisest ja saaks kuue võrrandi ja kuue tundmatuga süsteemi (eeldades, et on kuus takistit, millest kaks omavad lõpmata suurt väärtust), aga selle lahendamine on tüütu. Paneme tähele, et kui mõõta klemmide A ja B vahelist takistust ning lühistada klemmid B ja C, B ja D või C ja D vahel, siis oommeetri näit $R_1$ ei muutu. See on võimalik vaid siis, kui B, C ja D on kolmnurkühenduses, A ja B vahel on takistus RAB = $R_1$ ja ülejäänud takistused puuduvad. Mõõdame nüüd B ja C vahelise takistuse $r_1$ lühistades klemmid B ja D:

RB$-$1C + RC$-$$D_1$ = $r_1$$-$1;

analoogselt RC$-$$D_1$ + RD$-$1B = $r_2$$-$1 ja RD$-$1B + RB$-$1C = $r_3$$-$1. Summeerides leiame

RB$-$1C + RC$-$$D_1$ + RD$-$1B = 1 ($r_1$$-$1 + $r_2$$-$1 + $r_3$$-$1),

millest 1 RB$-$1C = ($r_1$$-$1 $-$ $r_2$$-$1 + $r_3$$-$1),

RC$-$$D_1$ = 1 ($r_1$$-$1 + $r_2$$-$1 $-$ $r_3$$-$1),

RD$-$1B = 1 ($-$$r_1$$-$1 + $r_2$$-$1 + $r_3$$-$1).
\probend
